% ---
% title: C++17 属性标记：fallthrough、nodiscard 与 maybe_unused
% date: 2026-07-13
% author: Crystal-Sky
% tags: C++, C++17, 基础语法, 属性, fallthrough, nodiscard, maybe_unused
% summary: 介绍 C++17 三个常用属性标记的作用、正确位置、典型场景、对比和常见错误。
% slug: cpp17-attributes
% course: C++
% course_slug: cpp
% chapter: 基础语法
% chapter_order: 1
% lesson_order: 7
% lesson_type: 基础语法
% ---

\documentclass[UTF8,12pt,a4paper]{ctexart}
\usepackage{geometry}
\usepackage{xcolor}
\usepackage{listings}
\usepackage{booktabs}
\usepackage{tcolorbox}
\geometry{left=2.5cm,right=2.5cm}

\lstset{
language=C++,
basicstyle=\ttfamily\small,
frame=single,
breaklines=true
}

\newtcolorbox{boxx}[1]{title=#1,colback=blue!3}

\title{C++17 属性标记\\
\texttt{[[fallthrough]]}、\texttt{[[nodiscard]]}、\texttt{[[maybe\_unused]]}}
\date{}

\begin{document}
\maketitle

\section{属性标记简介}

C++ 属性使用：

\begin{lstlisting}
[[attribute]]
\end{lstlisting}

向编译器提供额外信息。它通常不会改变程序逻辑，而是帮助编译器进行检查和优化。

C++17 常用属性：

\begin{itemize}
\item \texttt{[[fallthrough]]}
\item \texttt{[[nodiscard]]}
\item \texttt{[[maybe\_unused]]}
\end{itemize}

\section{\texttt{[[fallthrough]]}}

\subsection{作用}

用于 \texttt{switch} 中，表示当前 case 故意继续执行下一个 case。

普通情况：

\begin{lstlisting}
switch(x)
{
case 1:
    cout<<"one";

case 2:
    cout<<"two";
}
\end{lstlisting}

如果 x 为 1，会继续执行 case 2。

为了避免编译器认为忘记写 break：

\begin{lstlisting}
switch(x)
{
case 1:
    cout<<"one";
    [[fallthrough]];

case 2:
    cout<<"two";
    break;
}
\end{lstlisting}

\begin{boxx}{注意}
\texttt{[[fallthrough]]} 只是告诉编译器这里的穿透是有意的，不会改变执行流程。
\end{boxx}

\section{\texttt{[[nodiscard]]}}

\subsection{作用}

表示函数返回值或者对象结果不应该被忽略。

例如：

\begin{lstlisting}
[[nodiscard]]
int getValue()
{
    return 100;
}

int main()
{
    getValue(); // 编译器警告
}
\end{lstlisting}

常用于：

\begin{itemize}
\item 错误码；
\item 状态检查；
\item 资源创建函数。
\end{itemize}

例如：

\begin{lstlisting}
[[nodiscard]]
bool openFile();
\end{lstlisting}

调用者应该检查：

\begin{lstlisting}
if(openFile())
{
    // success
}
\end{lstlisting}

\subsection{修饰类型}

也可以修饰类型：

\begin{lstlisting}
struct [[nodiscard]] Error
{
    int code;
};
\end{lstlisting}

如果：

\begin{lstlisting}
Error getError();

getError();
\end{lstlisting}

同样会产生提醒。

\section{\texttt{[[maybe\_unused]]}}

\subsection{作用}

告诉编译器：

\begin{center}
这个变量、参数或者函数可能不会被使用。
\end{center}

例如：

\begin{lstlisting}
void debug([[maybe_unused]] int value)
{

}
\end{lstlisting}

不会产生 unused parameter 警告。

也可以用于变量：

\begin{lstlisting}
int main()
{
    [[maybe_unused]]
    int version = 17;
}
\end{lstlisting}

\section{三个属性对比}

\begin{center}
\begin{tabular}{lll}
\toprule
属性 & 用途 & 作用\\
\midrule
\texttt{fallthrough}
& switch
& 允许 case 穿透\\
\texttt{nodiscard}
& 返回值/类型
& 提醒不要忽略结果\\
\texttt{maybe\_unused}
& 变量/参数
& 允许未使用\\
\bottomrule
\end{tabular}
\end{center}

\section{综合示例}

\begin{lstlisting}
#include <iostream>
using namespace std;

[[nodiscard]]
bool connect()
{
    return true;
}

void test([[maybe_unused]] int mode)
{
}

int main()
{
    int x = 1;

    switch(x)
    {
    case 1:
        cout<<"start";
        [[fallthrough]];

    case 2:
        cout<<"run";
        break;
    }

    connect(); // 提醒返回值未使用

    return 0;
}
\end{lstlisting}

\section{常见错误}

\subsection{误解 nodiscard}

\texttt{nodiscard} 不会阻止程序运行，只产生编译器提醒。

\subsection{fallthrough 放置错误}

错误：

\begin{lstlisting}
case 1:
    [[fallthrough]];
    cout<<"hello";
\end{lstlisting}

正确：

\begin{lstlisting}
case 1:
    cout<<"hello";
    [[fallthrough]];

case 2:
\end{lstlisting}

\section{总结}

\begin{boxx}{快速记忆}

\begin{lstlisting}
[[fallthrough]]
// 我知道这里没有break

[[nodiscard]]
// 返回值不要丢

[[maybe_unused]]
// 这个变量可能不用
\end{lstlisting}

\end{boxx}

\end{document}
